\section{Applications to computer graphics}
For two-dimensional graphics, write a point as $(x,y,1)^T$. An affine
transformation is then represented by
\[
 H=\begin{bmatrix}a&b&h\\c&d&k\\0&0&1\end{bmatrix},\qquad
 H\begin{bmatrix}x\\y\\1\end{bmatrix}
 =\begin{bmatrix}ax+by+h\\cx+dy+k\\1\end{bmatrix}.
\]
The upper-left block controls the linear part and the last column controls
translation. A positive angle denotes counterclockwise rotation in the
usual oriented plane. To perform transformation $F$ first and $G$ second,
multiply their matrices as $GF$: the first action is the rightmost factor.
We use $R_\theta=\begin{bmatrix}\cos\theta&-\sin\theta\\
\sin\theta&\cos\theta\end{bmatrix}$ and write
$T_p=\begin{bmatrix}I_2&p\\0^T&1\end{bmatrix}$ for translation.

\practice{2.7.P1}{Rotate about a point other than the origin}
Let $p=(-2,6)^T$ and $\theta=-30^\circ$. Subtract $p$ to move the center
to zero, rotate, then add $p$ back. Algebraically this is
$x\mapsto R_\theta(x-p)+p=R_\theta x+(p-R_\theta p)$.
With $c=\sqrt3/2$, $s=-1/2$, the matrix is
\[
 T_p\begin{bmatrix}c&-s&0\\s&c&0\\0&0&1\end{bmatrix}T_{-p}
 =\boxed{\begin{bmatrix}
 \sqrt3/2&1/2&\sqrt3-5\\-1/2&\sqrt3/2&5-3\sqrt3\\0&0&1
 \end{bmatrix}}.
\]
For example $R_\theta p=(-\sqrt3+3,1+3\sqrt3)^T$, which explains the
translation column. Substituting $x=p$ gives $p$ again, as a rotation center
must be fixed.

\exercise{2.7.1}{Embed the shear into homogeneous coordinates}
Example 2 uses $(x,y)\mapsto(x+y/4,y)$. Preserve the final coordinate $1$:
\[
 \boxed{\begin{bmatrix}1&1/4&0\\0&1&0\\0&0&1\end{bmatrix}}.
\]
Multiplication by $(x,y,1)^T$ gives $(x+y/4,y,1)^T$, so the extra
coordinate introduces no translation and does not change the original shear.

\exercise{2.7.2}{Reflect and sketch a triangle}
Reflection through the $y$-axis negates $x$ and preserves $y$. Thus
\[
 \begin{bmatrix}-1&0\\0&1\end{bmatrix}
 \begin{bmatrix}4&2&5\\0&2&3\end{bmatrix}
 =\begin{bmatrix}-4&-2&-5\\0&2&3\end{bmatrix}.
\]
The original vertices $(4,0),(2,2),(5,3)$ therefore become
$(-4,0),(-2,2),(-5,3)$. Connect corresponding triples of vertices;
reflection preserves straight segments and lengths.
\begin{center}
\setlength{\unitlength}{5mm}
\begin{picture}(13,5)
\put(0,1){\vector(1,0){12}}\put(6,0){\vector(0,1){4.7}}
\put(6,1){\polyline(4,0)(2,2)(5,3)(4,0)}
\put(6,1){\polyline(-4,0)(-2,2)(-5,3)(-4,0)}
\put(10,0.3){\small$(4,0)$}\put(0.5,0.3){\small$(-4,0)$}
\put(8,3.2){\small$(2,2)$}\put(2.3,3.2){\small$(-2,2)$}
\put(10,4.3){\small$(5,3)$}\put(0,4.3){\small$(-5,3)$}
\put(12,0.5){$x$}\put(6.2,4.5){$y$}
\end{picture}
\end{center}

\exercise{2.7.3}{Translate first, rotate second}
Translation by $(2,1)$ takes $(x,y)$ to $(x+2,y+1)$.
A $90^\circ$ rotation then gives $(-(y+1),x+2)$.
Reading off coefficients yields
\[
 R_{90^\circ}T_{(2,1)}
 =\boxed{\begin{bmatrix}0&-1&-1\\1&0&2\\0&0&1\end{bmatrix}},
\]
where the rotation matrix is understood in its homogeneous $3\times3$ form.
The translated origin goes to $(-1,2)$, checking the last column.

\exercise{2.7.4}{Translate before unequal scaling}
The first action gives $(x-1,y+4)$, and the second gives
$((x-1)/2,3(y+4)/2)$. Therefore the matrix is
\[
 \boxed{\begin{bmatrix}1/2&0&-1/2\\0&3/2&6\\0&0&1\end{bmatrix}}.
\]
The translation column is scaled too, because the scaling is performed
\emph{after} translation.

\exercise{2.7.5}{Reflect, then rotate}
Reflection in the $x$-axis sends $(x,y)$ to $(x,-y)$.
Put $a=\sqrt2/2$. Rotation by $45^\circ$ gives
$(ax+ay,ax-ay)$, so
\[
 \boxed{\begin{bmatrix}a&a&0\\a&-a&0\\0&0&1\end{bmatrix}}.
\]
Equivalently its upper-left block is $R_{45^\circ}\operatorname{diag}(1,-1)$.

\exercise{2.7.6}{Reverse the preceding order}
First rotation gives $(ax-ay,ax+ay)$. Reflecting through the $x$-axis
now gives $(ax-ay,-ax-ay)$. Thus the matrix is
\[
 \boxed{\begin{bmatrix}a&-a&0\\-a&-a&0\\0&0&1\end{bmatrix}},\qquad a=\sqrt2/2.
\]
It differs from Exercise 5, illustrating why the order cannot be ignored.

\exercise{2.7.7}{Rotate sixty degrees around a fixed point}
Use the general formula $x\mapsto R_\theta x+p-R_\theta p$ with
$p=(6,8)^T$, $\theta=60^\circ$. Since
$R_\theta p=(3-4\sqrt3,3\sqrt3+4)^T$, the result is
\[
 \boxed{\begin{bmatrix}
 1/2&-\sqrt3/2&3+4\sqrt3\\
 \sqrt3/2&1/2&4-3\sqrt3\\0&0&1
 \end{bmatrix}}.
\]
Its action fixes $(6,8)$, which independently checks the translation terms.

\exercise{2.7.8}{Rotate forty-five degrees around another point}
Now $p=(3,7)^T$, and $R_{45^\circ}p=(-2\sqrt2,5\sqrt2)^T$.
Therefore
\[
 \boxed{\begin{bmatrix}
 \sqrt2/2&-\sqrt2/2&3+2\sqrt2\\
 \sqrt2/2&\sqrt2/2&7-5\sqrt2\\0&0&1
 \end{bmatrix}}.
\]
This is $T_pR_{45^\circ}T_{-p}$, not $R_{45^\circ}T_p$.

\exercise{2.7.9}{Precompute a composite transformation}
For $A(BD)$, the $2\times100$ product $BD$ has $200$ entries,
each requiring two multiplications: $400$ multiplications. Multiplication
by $A$ costs another $400$, for $800$ total.
For $(AB)D$, forming the four entries of $AB$ costs $4\cdot2=8$,
and multiplying its result by $D$ costs $400$, for $408$ total.
Thus precomputing the $2\times2$ composite saves $392$ multiplications.
The advantage grows when the same transformations are applied to more
vertices: combine the small transformation matrices once, then transform
all data points with their product.

\exercise{2.7.10}{Which geometric transformations commute?}
Let $D(x)=\lambda x$, $R(x)=Rx$, and $T(x)=x+p$.
Then $D(Rx)=\lambda Rx=R(\lambda x)$, so a uniform dilation and an
origin-centered rotation always commute.
By contrast,
\[
 D(Tx)=\lambda x+\lambda p,\qquad T(Dx)=\lambda x+p.
\]
They commute exactly when $(\lambda-1)p=0$, and generally do not.
Likewise $R(Tx)=Rx+Rp$ and $T(Rx)=Rx+p$, so rotation and translation
commute exactly when $Rp=p$. A nontrivial planar rotation fixes only the
zero vector, making a nonzero translation fail this test.
For concrete failures take $p=(1,0)$, $\lambda=2$, and a $90^\circ$
rotation, and evaluate the two orders at $x=0$.

\exercise{2.7.11}{Two shear-and-scale factors produce a rotation}
The first action is $A_1$, the second is $A_2$, so calculate $A_2A_1$.
Its upper-left block is
\[
 \begin{bmatrix}\sec\varphi&-\tan\varphi\\0&1\end{bmatrix}
 \begin{bmatrix}1&0\\\sin\varphi&\cos\varphi\end{bmatrix}
 =\begin{bmatrix}
 \sec\varphi-\tan\varphi\sin\varphi&-\tan\varphi\cos\varphi\\
 \sin\varphi&\cos\varphi
 \end{bmatrix}
 =\begin{bmatrix}\cos\varphi&-\sin\varphi\\\sin\varphi&\cos\varphi\end{bmatrix}.
\]
The first simplification uses
$(1-\sin^2\varphi)/\cos\varphi=\cos\varphi$.
The last homogeneous row and column remain $(0,0,1)$, so this is rotation
by $\varphi$. The stated factors require $\cos\varphi\ne0$; the rotation
itself still exists at other angles, but these particular factors do not.

\exercise{2.7.12}{Three shears can also implement rotation}
Write $t=\tan(\varphi/2)$, $s=\sin\varphi$. Ignoring the unchanged
homogeneous coordinate for a moment, multiply the three $2\times2$ blocks:
\[
 \begin{bmatrix}1&-t\\0&1\end{bmatrix}
 \begin{bmatrix}1&0\\s&1\end{bmatrix}
 \begin{bmatrix}1&-t\\0&1\end{bmatrix}
 =\begin{bmatrix}1-ts&-2t+t^2s\\s&1-st\end{bmatrix}.
\]
Half-angle identities give $st=1-\cos\varphi$ and
$-2t+t^2s=-\sin\varphi$, so the result is $R_\varphi$.
For instance $\sin\varphi=2t/(1+t^2)$ makes the second identity
$-2t+2t^3/(1+t^2)=-2t/(1+t^2)$.
Appending the homogeneous row and column yields precisely the claimed
$3\times3$ rotation matrix. Each shear needs one nontrivial multiplication
per point, whereas the usual rotation needs four. The factorization requires
$\tan(\varphi/2)$ to be defined, excluding odd multiples of $180^\circ$.

\exercise{2.7.13}{Factor a general affine transformation}
Block multiplication gives
\[
 \begin{bmatrix}A&p\\0^T&1\end{bmatrix}
 =\begin{bmatrix}I_2&p\\0^T&1\end{bmatrix}
  \begin{bmatrix}A&0\\0^T&1\end{bmatrix}.
\]
The right factor acts first and sends $(x,1)$ to $(Ax,1)$.
The left then adds $p$, giving $(Ax+p,1)$. Thus the transformation is
a linear transformation of the plane followed by a translation, as requested.
No inverse or invertibility of $A$ is needed.

\exercise{2.7.14}{Separate the linear rotation from its final translation}
The transformation from Exercise 7 is
$x\mapsto R_{60^\circ}x+(3+4\sqrt3,4-3\sqrt3)^T$.
Hence it is rotation about zero followed by translation by
\[
 \boxed{p=(3+4\sqrt3,\ 4-3\sqrt3)^T}.
\]
The translation is the original center minus its rotated image; it is
not simply the center $(6,8)$ itself.

\exercise{2.7.15}{Normalize a homogeneous four-vector}
A homogeneous point $(X,Y,Z,w)$ with $w\ne0$ represents
$(X/w,Y/w,Z/w)$. Dividing by $1/24$ multiplies the first three entries
by $24$, so the represented vector is
$\boxed{(12,-6,-3)^T}$ in $\R^3$.

\exercise{2.7.16}{Different representatives of one point}
Yes: $(10,-20,-30,40)=10(1,-2,-3,4)$.
Dividing either representative by its final coordinate gives
$(1/4,-1/2,-3/4)$. A common nonzero scale changes the homogeneous
representative but not the point it represents.

\exercise{2.7.17}{A three-dimensional rotation about the x-axis}
The $x$-coordinate is unchanged, while the $(y,z)$ coordinates undergo
the usual $60^\circ$ planar rotation in the orientation shown. The
homogeneous matrix is
\[
 \boxed{\begin{bmatrix}
 1&0&0&0\\0&1/2&-\sqrt3/2&0\\
 0&\sqrt3/2&1/2&0\\0&0&0&1
 \end{bmatrix}}.
\]
It sends $e_2$ toward $e_3$, confirming the positive orientation about
$e_1$ by the right-hand convention.

\exercise{2.7.18}{Rotate about z, then translate}
The rotation leaves $z$ fixed and has planar angle $-30^\circ$.
Translation after rotation appends $(5,-2,1)^T$ as the last column:
\[
 \boxed{\begin{bmatrix}
 \sqrt3/2&1/2&0&5\\-1/2&\sqrt3/2&0&-2\\0&0&1&1\\0&0&0&1
 \end{bmatrix}}.
\]
For example the origin is sent to $(5,-2,1)$, checking the order of actions.

\exercise{2.7.19}{Project a triangle from a finite viewing point}
A line from $c=(0,0,10)$ through $v=(x,y,z)$ is
$c+t(v-c)=(tx,ty,10+t(z-10))$.
Its intersection with the image plane $z=0$ has $t=10/(10-z)$.
Consequently the projection formula is
\[
 (x,y,z)\longmapsto\left(\frac{10x}{10-z},\frac{10y}{10-z},0\right).
\]
For the three given vertices the scale factors are $5/3,5/4,5/2$,
so their images are
\[
 \boxed{(7,2,0),\quad(15/2,5,0),\quad(5,5,0)}.
\]
All original points have $z<10$, so each line meets the image plane,
and the triangular region projects to the triangle with these vertices.

\exercise{2.7.20}{A second perspective triangle}
Apply the same formula. At $(7,3,-5)$ the scale factor is $2/3$;
at $(12,8,2)$ it is $5/4$; at $(1,2,1)$ it is $10/9$.
Thus the image triangle has vertices
\[
 \boxed{(14/3,2,0),\quad(15,10,0),\quad(10/9,20/9,0)}.
\]
Perspective projection is not a single ordinary $3\times3$ linear map
on the original three coordinates: the required division depends on $z$.
Homogeneous coordinates encode that division by a final normalization.

\exercise{2.7.21}{Convert standard color data into screen data}
Let $M$ denote the stated CIE conversion matrix, so
$M(R,G,B)^T=(X,Y,Z)^T$. To recover the gun data solve this system:
\[
 \boxed{\begin{bmatrix}R\\G\\B\end{bmatrix}
 =M^{-1}\begin{bmatrix}X\\Y\\Z\end{bmatrix}},\qquad
 M^{-1}=\frac1{20224}
 \begin{bmatrix}45677&-21023&-7023\\-27293&47407&1407\\1840&-6160&25840\end{bmatrix}.
\]
This exact rational inverse treats the printed decimals as exact input.
Its decimal approximation is
\[
 \begin{bmatrix}2.258554&-1.039508&-.347261\\
 -1.349535&2.344096&.069571\\.090981&-.304589&1.277690\end{bmatrix}.
\]
Exact matrix multiplication checks $MM^{-1}=I$; in particular the
conversion is uniquely reversible. A matrix program can reproduce it with
SymPy rational decimal entries and \texttt{M.inv()}; do not reverse the
conversion by multiplying by $M$ a second time.

\exercise{2.7.22}{Recover RGB from a television signal}
Write the supplied matrix as $N$, so the transmitted vector is
$(Y,I,Q)^T=N(R,G,B)^T$. Solving for RGB gives
\[
 \boxed{\begin{bmatrix}R\\G\\B\end{bmatrix}
 =N^{-1}\begin{bmatrix}Y\\I\\Q\end{bmatrix}},
\]
\[
 N^{-1}=\frac1{50845523}
 \begin{bmatrix}
 51002600&48549800&31415400\\
 50681600&-13764200&-32784600\\
 51277600&-56463200&86415400
 \end{bmatrix}
 \approx\begin{bmatrix}
 1.003089&.954849&.617860\\
 .996776&-.270706&-.644788\\
 1.008498&-1.110485&1.699568
 \end{bmatrix}.
\]
The exact form was calculated and verified by rational matrix multiplication.
It uses the coefficients as printed, including $.212,-.528,.311$ in the
last row. Different engineering conventions or more accurate calibration
coefficients would produce a different numerical inverse; no such change
has been silently substituted here.
